\documentclass[11pt]{article}
\usepackage[T1]{fontenc}
\usepackage{lmodern}
\usepackage[margin=1in]{geometry}
\usepackage{amsmath,amssymb,amsthm}
\usepackage[hidelinks]{hyperref}
\newtheorem{definition}{Definition}[section]
\newtheorem{proposition}[definition]{Proposition}
\newtheorem{theorem}[definition]{Theorem}
\newtheorem{remark}[definition]{Remark}
\newcommand{\Cat}[1]{\mathsf{#1}}
\newcommand{\R}{\mathbb R}
\newcommand{\N}{\mathbb N}
\newcommand{\Ty}{\operatorname{Ty}}
\newcommand{\Tm}{\operatorname{Tm}}
\newcommand{\El}{\operatorname{El}}
\newcommand{\norm}[1]{\left\lVert #1\right\rVert}
\newcommand{\ev}{\operatorname{ev}}
\setlength{\emergencystretch}{2em}
\title{A Task Specific Elaboration Situation for Changing Model Languages}
\author{}
\date{4 November 2026}
\begin{document}
\maketitle

\begin{abstract}
The problem is to turn a proposed change from a bounded linear weighted
finite automaton class to a rational recurrence into a checkable mathematical
and evidential claim. An interface for the whole relevant elaborator is organized as a diagram
of displayed categories: verification contexts and dependent syntax, task
data and model denotations, claims and proofs, dependency telescopes and typed
holes, transformations, and optional canonicity witnesses all occur within
one diagram. A conditional construction from supplied task-record presheaves and a
category with families makes substitution and hole-filling laws precise. A two-by-two Hankel example
then supplies an end-to-end comparison with a rational certificate and
an explicit residual empirical assumption. The record signatures and checker interface are schematic where a concrete
encoding has not been supplied. The categorical template and the numerical
comparison have proved fragments; this is not an implemented elaborator
or a theorem that every useful representation change has this shape.
\end{abstract}

\section{The actual claim to be elaborated}

A linear Hankel rank obstruction is relative to a linear realization class.
The behavior
\[
 g(a^n)=\frac1{n+1},\qquad x_0=1,\qquad
 x_{n+1}=\frac{x_n}{1+x_n},\qquad g(a^n)=x_n
\]
has a one-coordinate rational realization on $x>0$, but no finite real
linear WFA realization. Thus a change to this recurrence is not a
behavior-preserving compilation of the target into a finite WFA. A useful
elaborator must distinguish:
\begin{enumerate}
\item an exact interpretation or coordinate change preserving a stated
      denotation;
\item a model replacement supported by a comparison on specified evidence;
\item an identification of the true mechanism or a canonical realization.
\end{enumerate}
The first two have concrete constructions below. Neither entails the third.

The proposed intermediate representation retains exactly the distinctions
needed to prevent those conflations. Dependency graphs, verification
judgments, and contextual syntax are not substitutes for the categorical
diagram. They are the data represented by its displayed categories.

\section{A diagram schema for this elaboration situation}

\subsection{Verification theory and contextual syntax}

For the categorical template, assume a verification theory $T$ whose
chosen contextual representation is a split category with families
\[
 (\Cat C_T,\Ty_T,\Tm_T,1,(-).(-),p,q).
\]
An object $\Gamma$ of $\Cat C_T$ is a context. When this is the
syntactic category of $T$, it is a well-formed syntactic context. An arrow
$\sigma:\Delta\to\Gamma$ is a typed substitution, so
\[
 A[\sigma]\in\Ty_T(\Delta),\qquad
 t[\sigma]\in\Tm_T(\Delta,A[\sigma])
 \quad\text{if }t\in\Tm_T(\Gamma,A).
\]
Substitution preserves identities and composition. Context extension
$\Gamma.A$ includes its projection and generic variable, with the usual
comprehension universal property. This is where dependent hypotheses and context extension live. Binders
require the corresponding type formers and their substitution laws;
context extension alone does not supply them. The mathematical fragment requires natural numbers,
real arithmetic, finite matrices, sequences, norms, propositions,
quantification, and the relevant proved linear-algebra lemmas.

An abstract semantic category with families does not by itself provide
a grammar, a typing derivation system, or an executable checker. A
syntactic implementation needs those data, including a specified checker $V_T$ for the
judgment $\Gamma\vdash t:A$, a semantic interpretation of the theory,
and its soundness theorem. The template below assumes the chosen contextual laws; it does not
construct or verify a syntactic category for a particular proof assistant. Merely drawing its nodes does not build
that checker or prove its soundness. Checker soundness means that an
accepted judgment is derivable in $T$; semantic soundness of $T$ is a
separate assertion about its interpretation. A small Lean-checked fragment is
provided separately; it does not certify the entire diagram.

For any presheaf $F:\Cat C_T^{\mathrm{op}}\to\Cat{Set}$, use its
category of elements $\El(F)$. An object is $(\Gamma,x)$ with
$x\in F(\Gamma)$; a morphism over $\sigma:\Delta\to\Gamma$ exists
exactly when $x'=F(\sigma)(x)$. Its identity and composition are inherited
from substitution. This gives an explicit displayed category, not just a
name for a box. Dependent record fields are formed by iterating this
construction over the appropriate total category.

\subsection{The task tower}

The following task-record list is a signature schema. For the categorical
result, \emph{assume supplied} a presheaf of data records, a dependent
presheaf of model records over it, a dependent presheaf of goals, and a
substitution-stable family of theory-admitted proof terms, including
any requested canonicity witnesses. Transformation records and their
packaging maps must also be supplied. All displayed
record fields and maps must be natural under the chosen substitution.
Once these data are supplied, iterated categories of elements give:
\[
\begin{array}{ccccccccc}
 \Cat{Can} &\longrightarrow& \Cat{Prf} &\longrightarrow&
 \Cat{Goal} &\longrightarrow& \Cat{Rep} &\longrightarrow& \Cat{Data}\\
 &&&&&&&&\downarrow\\[-2pt]
 &&&&&&&&\Cat C_T .
\end{array}                                                     \tag{D1}
\]
Each horizontal projection forgets the final dependent record field.
The field descriptions below do not alone prove that a concrete encoding
exists or that its operations are natural. Those are explicit implementation
obligations. The intended fibers have the following content.
\begin{description}
\item[$\Cat{Data}$] A finite observation context, an observation matrix,
a norm and normalization, an uncertainty specification, the competitor
class and its resource budget, and a stakeholder objective. A latent true
behavior can be a variable. Noise bounds are propositions in context,
not truths manufactured from the observation matrix.
\item[$\Cat{Rep}$] A finite ordered tuple of model records, each carrying a language tag,
a well-typed presentation, its state domain and initial condition, and a denotation in the behavior
space. Rational denominators require domain obligations. Candidate
semantics must factor through an actual model; an arbitrary low-rank
matrix is not a realization witness.
\item[$\Cat{Goal}$] A proposition with all these data in scope. Typical
fibers are robust class-relative comparison, exact interpretation,
approximate interpretation with a stated error budget, or realization
uniqueness under a specified equivalence relation.
\item[$\Cat{Prf}$] A theory-admitted term at the preceding proposition; in the syntactic
case this means a typing derivation. Checker acceptance is separate. An
assumption is represented by a variable in the context; it is not silently
converted into an unconditional theorem. Every conclusion retains the
context on which its proof depends.
\item[$\Cat{Can}$] An optional proof of a \emph{stated} canonicity property,
such as uniqueness up to a specified isomorphism among admissible outputs.
Its fiber can be empty. An unfilled witness is not a proof that no witness
exists. The projection never creates this proof merely because an
elaboration succeeded.
\end{description}
Transformations and interpreters form another displayed branch:
\[
\begin{array}{ccccc}
 \Cat{TransPrf}&\longrightarrow&\Cat{TransGoal}
   &\longrightarrow&\Cat{Rep}\mathbin{\times_{\Cat{Data}}}\Cat{Rep}\\
 \downarrow&&\downarrow&&\downarrow\\
 \Cat{Prf}&\longrightarrow&\Cat{Goal}&\longrightarrow&\Cat{Data}.
\end{array}                                                     \tag{D2}
\]
A transformation record gives the input and output languages/models,
its map or interpreter, its domain, and its exact semantic equation or
approximation bound. For the vertical map to $\Cat{Goal}$, concatenate the two ordered tuples
of endpoint records and retain the transformation equation as the goal.
Thus the underlying $\Cat{Rep}$ record of that ordinary goal is exactly
the ordered concatenation, which is natural under componentwise substitution.
This packaging, and the map forgetting a transformation certificate to
its underlying proof, are required to commute with substitution. Thus the
left square is the forgetful square for that declared proposition; it
is not an unspecified choice of one endpoint.
Composition is available only under the conditions in Section~\ref{sec:transport}.

\subsection{Telescopes, holes, and solutions are in the same diagram}

Let $\Cat{Tel}_T$ be the category of labeled telescope projections
$\pi_\Delta:\Gamma.\Delta\to\Gamma$, where
$\Delta=(x_1:A_1,\ldots,x_k:A_k)$ is a finite dependent telescope over
$\Gamma$. A morphism from $(\Gamma',\Delta')$ to $(\Gamma,\Delta)$
is a pair
\[
 \sigma:\Gamma'\to\Gamma,\qquad
 u:\Gamma'.\Delta'\to\Gamma.\Delta,
 \qquad \pi_\Delta u=\sigma\pi_{\Delta'}.
\]
Identities and composition are componentwise. Equivalently this is the
category of the selected labeled objects of the arrow category of
$\Cat C_T$, with all commuting squares between them. Define functors
\[
 p(\Gamma,\Delta)=\Gamma,\quad p(\sigma,u)=\sigma,\qquad
 e(\Gamma,\Delta)=\Gamma.\Delta,\quad e(\sigma,u)=u.
\]
Thus $p$ forgets the local telescope and $e$ extends the context.
Pulling the judgment and proof categories back along $e$ yields the
open-goal and open-proof vertices. A solution is a section
$s:\Gamma\to\Gamma.\Delta$ of the telescope projection.
Define $\Cat{Sol}$ to have these section-equipped objects, with a
morphism above $(\sigma,u)$ precisely when
\[u s'=s\sigma.\]
This condition is preserved by identities and composition. The projection
$\Cat{Sol}\to\Cat{Tel}_T$ is a displayed category, but is not
asserted to be a discrete fibration or to have a lift over every square.
The split category with families supplies reindexing along its chosen
cartesian squares only. These definitions give
three pullback constructions:
\[
\begin{array}{ccccc}
 \Cat{OpenPrf}&\longrightarrow&\Cat{Prf}
 &\qquad&\Cat{Filled}\longrightarrow\Cat{Sol}\\
 \downarrow&&\downarrow&&\downarrow\qquad\qquad\downarrow\\
 \Cat{Tel}_T&\xrightarrow{\ e\ }&\Cat C_T
 &&\Cat{OpenPrf}\longrightarrow\Cat{Tel}_T .
\end{array}                                                     \tag{D3}
\]
The analogous left-hand square with $\Cat{OpenGoal}$ and $\Cat{Goal}$
defines the open goals. Explicitly,
\begin{align*}
\Cat{OpenPrf}&=\Cat{Tel}_T\times_{e,\Cat C_T}\Cat{Prf},\\
\Cat{OpenGoal}&=\Cat{Tel}_T\times_{e,\Cat C_T}\Cat{Goal},\\
\Cat{Filled}&=\Cat{OpenPrf}\times_{\Cat{Tel}_T}\Cat{Sol}.
\end{align*}
The filling functor sends a morphism above $(\sigma,u)$ to substitution
by $\sigma$. It is a further arrow in the same diagram:
\[
\begin{array}{ccc}
 \Cat{Filled}&\xrightarrow{\ \mathrm{fill}\ }&\Cat{Prf}\\
 \downarrow&&\downarrow\\
 \Cat{Tel}_T&\xrightarrow{\ p\ }&\Cat C_T,
\end{array}\qquad
 (\Gamma,\Delta,A,t,s)\longmapsto(\Gamma,A[s],t[s]).             \tag{1}
\]
Equation (1) abbreviates the full task record: every data, model, and
goal field is substituted by $s$, not only the displayed type and term.
The proof projection in (D3) has context $e(\Gamma,\Delta)$;
the filled proof in (1) has context $p(\Gamma,\Delta)$.
They are related by the natural section $s:p\to e$ after both functors
are pulled back to solution records, not by equating these different contexts. Thus a dependent output type
is not treated as fixed before its holes are solved.

The vertical arrows of $\Cat{Tel}_T$ above $\mathrm{id}_\Gamma$
are precisely the partial-solution substitutions. Thus a partial solution
to $\Delta$ is a substitution
$s:\Gamma.\Delta'\to\Gamma.\Delta$ over $\Gamma$. Composition of
partial solutions is substitution composition. The operational step
from an old obligation to a residual one has the opposite direction to
this semantic substitution. A dependency graph is an acyclic presentation
of the telescope: an edge records actual free-variable dependence.
Different topological orders are interchangeable only by the appropriate
exchange substitutions; arbitrary dependent declarations cannot be reordered.
A cycle needs an explicitly justified recursive or fixed-point type,
not a topological sort.

\begin{proposition}[Conditional categorical template and filling]
Assume the split category with families and the substitution-stable task
presheaves and natural maps specified above have been supplied. Construct
$\Cat{Tel}_T$ and $\Cat{Sol}$ by the commuting-square and section
constructions. Then (D1)--(D3) are diagrams of categories and functors,
and the filling map (1) is a functor. A theory-admitted open term $t:A$
in $\Gamma.\Delta$ and a section $s$ give a theory-admitted term
$t[s]:A[s]$ in $\Gamma$. Staged filling agrees with composite filling.
If the theory has a sound semantic interpretation, the resulting
judgment is valid under its remaining assumptions. If an executable
checker is additionally to accept the filled term, its acceptance must
be proved closed under this substitution operation or its completeness
for these derivations must be supplied; checker soundness alone is
insufficient.
\end{proposition}
\begin{proof}
Categories of elements construct the task tower from the supplied
presheaves. Commuting squares give the category laws for
$\Cat{Tel}_T$. The equation $us'=s\sigma$ is preserved by composition,
so it defines the section category. Ordinary categorical pullbacks give
the open and filled vertices. A morphism of open proofs has
$t'=t[u]$ and $A'=A[u]$. Hence
\[
 t'[s']=t[u][s']=t[us']=t[s\sigma]=t[s][\sigma],
\]
and likewise for types and all dependent record fields. This proves
that filling is a functor. The substitution rule proves the filled
judgment, and its composition law proves staged filling. Semantic
soundness applies without erasing the remaining context. The statement
about actual checker acceptance explicitly requires its additional
interface, rather than deriving completeness from soundness.
\end{proof}

\paragraph{Elaboration control and proof status.}
The elaborator is a possibly branching search through these typed partial
solutions, not a presumed total function from prose to a canonical answer.
An untrusted proposal becomes a certified step only after its substitution
and proof obligations pass $V_T$. Operational paths run from an old obligation to a residual one; their
functor therefore targets the \emph{opposite} of the semantic
partial-solution category above. It does not reverse the substitution
arrows inside that semantic category. Failed proposals are retained as failed
attempts, not as morphisms. Proof status has its own displayed audit branch over the goal and
partial-solution records. A fiber contains the exact source judgment,
its open fields and assumptions, the proposed or accepted proof artifact,
the checker and library version when applicable, and the status
``open'', ``proved on paper'', or ``kernel checked''.
\[
 \Cat{Audit}\longrightarrow\Cat{Goal},\qquad
 \Cat{Check}\longrightarrow\Cat{ProposedStep}\longrightarrow\Cat{OpenGoal}.
                                                               \tag{D4}
\]
All arrows in (D4) use the semantic substitution direction. Operational
logs use the opposite categories over the opposite base categories, just
as the control paths above do. Untrusted proposals are objects carrying
strings and goals, not automatically base-category morphisms.
An accepted-check fiber contains a checker acceptance witness for that
exact proposal; other fibers can be empty. An audit morphism must supply
the valid proof transport or replay record for the changed judgment.
These displayed categories are not asserted to be discrete fibrations:
there need not be an audit lift over every proposed substitution.
Identities retain the artifact and composition concatenates valid
transport records. The soundness interface connects accepted artifacts
to $\Cat{Prf}$; no arrow upgrades a paper proof or an empirical claim
solely by changing its status label. These provenance fields are outside
the proposition being proved, but inside the elaboration diagram.
Formally, first supply a graph of accepted primitive records with exact
source and target judgments and a valid semantic-direction base morphism
label on each edge. The free path category then has the induced functor
to that judgment category. For operational-direction edges, use the
opposite base instead. A morphism above a substitution is a path whose substitutions
compose to it. Empty paths and concatenation give the category laws.
This constructs the audit category relative to an acceptance relation;
it does not implement or prove sound that relation.

The diagram is an interface schema for the whole \emph{relevant} ad hoc
elaborator, including its verification theory, syntax, semantics,
dependencies, search state, and output certificates. The exact theorem is
about the categorical template \emph{when its stated data are supplied}.
It does not establish that the full task signature has been encoded as
those presheaves in an implemented theory, or supply a parser, tactic
scheduler, or decision procedure for its open goals. Across verification
theories it is a family of such diagrams. Transport between two family
members requires an actual interpretation preserving contexts,
substitution, judgments, and the required semantics; none is assumed here.

\section{Models and observations}

Fix a finite alphabet $\Sigma$. Put $\Cat B=\R^{\Sigma^*}$.
For finite prefix and suffix sets $c=(S,E)$ define
\[
 B_c=\R^{S\times E},\qquad
 (\ev_c f)_{s,e}=f(se).
\]
The repeated-word equalities are part of the image of $\ev_c$.
An arbitrary matrix in $B_c$ need not lie in that image, much less arise
from a transition-consistent WFA.

The WFA language of budget $r$ denotes tuples
$(\alpha,(A_a),\beta)$ of dimension at most $r$ with
$f(w)=\alpha A_w\beta$. The rational-recurrence language denotes a
state domain $D$, an initial state in $D$, rational update maps $R_a$,
a rational readout $o$, and proofs that denominators do not vanish on
$D$ and that $R_a(D)\subseteq D$. Its word behavior is defined by
iteration. Presentations without these proofs remain in an open telescope.
The proof obligations can be difficult or undecidable; the language tag
is not an automatic discharge mechanism.

For $c\subseteq d$, coordinate restriction
$\rho_{dc}:B_d\to B_c$ satisfies
$\rho_{dc}\ev_d=\ev_c$. For nonnegative radii $\epsilon_c,\epsilon_d$ and a nonnegative
certified Lipschitz constant $K$,
\[
 \norm{\rho_{dc}Y_d-Y_c}+K\epsilon_d\leq\epsilon_c             \tag{2}
\]
suffices to transport the uncertainty ball at $d$ into that at $c$.
For operator norm on rectangular submatrices $K=1$ is valid.
For $m\times m$ matrices with normalized Frobenius norm
$\norm{X}_{\nu,m}=\norm{X}_F/m$, restriction from size $m'$ to size
$m$ has bound $K=m'/m$, not generally $1$.
Morphisms carrying $(\rho,K)$ compose with constants multiplied.
The triangle inequality proves that certificates (2) compose.

The category of evidence states has objects $(c,Y,\epsilon)$ with
$\epsilon\geq0$. An arrow from $(d,Y_d,\epsilon_d)$ to $(c,Y_c,\epsilon_c)$ is
$(\rho_{dc},K)$, where $K\geq0$ is a valid Lipschitz bound for the
stated restriction and (2) holds, with a fixed norm at each context. Treat certificate
validity here as a property, rather than keeping arbitrary proof terms
as unequal parallel arrow data. The identity has $K=1$; composition
multiplies constants, and associativity is ordinary multiplication
associativity. Proof artifacts can be retained separately in the audit
path category. This category need not have an arrow for every proposed update. In particular,
small residual on old observations supplies no certificate on newly
added observations. New evidence is an extension of the task data,
followed by a new proof obligation.

\section{Exact and approximate transformations}
\label{sec:transport}

\begin{definition}[Exact semantic interpretation]
For model languages $L,K$ with denotations $d_L,d_K$, a certified
interpretation consists of a map $F:M_L\to M_K$ and a proof
$d_K(F(m))=d_L(m)$ for each allowed $m$. Composition is function
composition with the composed equality proof. Restricted-domain
interpretations state and prove the domain and image conditions.
\end{definition}

For state presentations a sufficient witness is a map $h:D\to D'$
with
\[
 h(x_0)=x'_0,\qquad hR_a=R'_a h,\qquad o=o'h.
\]
Induction on words proves behavioral preservation. If $h$ is invertible
with these properties it is a presentation isomorphism. For the example,
$h(x)=1/x$ is an isomorphism on the positive domain between
$x'=x/(1+x),\ o(x)=x$ and $z'=z+1,\ o'(z)=1/z$.
It preserves the dynamics and readout; it does not turn the readout into
a linear one. A WFA does embed into a same-dimensional rational state
language using linear updates and readout, but this has no general inverse.

At a fixed observation context, an approximate interpretation may instead
carry
\[\norm{\ev_c d_K(F(m))-\ev_c d_L(m)}\leq\delta.\]
Two such interpretations compose with error at most
$\delta_1+\delta_2$ in the same observation norm, provided the first
map lands in the second map's admissible domain and the second bound
holds at that actual intermediate model. A uniform bound on that domain
is sufficient; a bound for an unrelated input is not. An observation map with Lipschitz
constant $K$ multiplies that budget by $K$. Approximate agreement within
a fixed tolerance is not transitive and must not be silently quotiented
as an equivalence relation.

A comparison certificate does not supply any of these maps. If it asserts
that a candidate $g$ has lower loss than every model of a bounded class,
it is a proof in $\Cat{Goal}$ relating two denotations on the given task.
Selecting $g$ is a new model choice, not a semantic equality between $g$
and the previous model. A further valid rule is a strict language exit:
if $g$ is not finite-WFA realizable, $h$ is WFA realizable, and $t\ne0$,
then $(1-t)h+tg$ is not finite-WFA realizable. Otherwise linear closure
would realize $g=((1-t)h+tg-(1-t)h)/t$. A rational product-state
construction can realize the mixture when both components are rational
recurrences. This rule preserves neither the old behavior nor its
language; its proper type records the changed output and the exit proof.

\paragraph{Compatibility with existing capacity operations.}
The original WFA note's gated direct sum, evidence-image span, and
reverse-mode pullback have different types in this diagram. A gated
extension is a parameterized model construction with a specified equality
at initialization. The observational span retains generators and tests;
it is not an old-state embedding. A differentiable parameter map
$i:\Theta\to\Theta'$ transports covectors by $Di^*$ when its semantic
square commutes. It does not imply optimizer commutation. These records
can coexist, but a successful proof of one cannot fill a hole of another
type.

\section{The comparison rule and its dependency telescope}

Let $m\geq1$ and $0\leq r<m$. Let $Y$ be an $m\times m$ observed Hankel matrix and use
$\norm{X}_\nu=\norm{X}_F/m$. Let $f_*$ be the unknown target and
$G=\ev_c g$ the block of an actual candidate model. Define
\[
 b_r(Y)=\frac1m\left(\sum_{j>r}\sigma_j(Y)^2\right)^{1/2}.
\]
The elaborator uses certified bounds $\beta\leq b_r(Y)$ and
$c_g\geq\norm{Y-G}_\nu$, rather than treating a floating-point SVD
as a proof. The essential telescope is
\begin{align*}
 &c,m,r,Y,f_*,g,\eta,\beta,c_g;\\
 &h_{\mathrm{model}}:g\text{ has the declared well-defined denotation};\\
 &h_{\mathrm{noise}}:\norm{Y-\ev_c f_*}_\nu\leq\eta;\\
 &h_{\mathrm{spectral}}:\beta\leq b_r(Y);\qquad
 h_{\mathrm{fit}}:\norm{Y-G}_\nu\leq c_g;\\
 &h_{\mathrm{nonnegative}}:0\leq\eta\ \wedge\ 0\leq c_g;\qquad
 h_{\mathrm{margin}}:c_g+2\eta<\beta .
\end{align*}
All later fields depend on their preceding parameters. The mathematical
rule also consumes proved WFA factorization and low-rank approximation
lemmas from the verification theory.

\begin{theorem}[Conditional class-relative robust replacement]
Under these premises every WFA $A$ of dimension at most $r$ satisfies
\[
 L_{f_*}(A)-L_{f_*}(g)\geq
 \Gamma:=(\beta-\eta)^2-(c_g+\eta)^2>0,
 \qquad L_{f_*}(h)=\norm{\ev_c h-\ev_c f_*}_\nu^2.
                                                               \tag{3}
\]
\end{theorem}
\begin{proof}
The WFA block has rank at most $r$. The best rank-$r$ approximation
bound gives
\[\norm{Y-\ev_c A}_\nu\geq b_r(Y)\geq\beta.\]
The triangle inequality implies
\[\norm{\ev_c f_*-\ev_c A}_\nu\geq\beta-\eta,\qquad
\norm{\ev_c f_*-G}_\nu\leq c_g+\eta.\]
The margin makes both comparison bounds nonnegative, with the former
strictly larger. Squaring and subtracting proves (3).
\end{proof}

No independence between $g$ and $Y$ was used: this deterministic theorem
holds simultaneously for any fitted candidate satisfying its premises.
A probabilistic claim still needs a justified event containing
$h_{\mathrm{noise}}$ and its probability. Reusing data for candidate
selection does not by itself invalidate the deterministic conclusion,
but a pointwise statistical noise guarantee cannot silently be promoted
to a simultaneous or optional-stopping guarantee.

The comparison is relative to the norm, horizon, and WFA budget in its
type. Changing any of them creates a new goal. It does not prove global
superiority, minimal physical state dimension, or that a rational language
was the only sensible enlargement.

\section{An end to end finite check}

Take the unary alphabet, $S=E=\{\varepsilon,a\}$, $m=2$, $r=1$, and
\[
 Y=G=\begin{pmatrix}1&1/2\\1/2&1/3\end{pmatrix},
 \qquad g(a^n)=1/(n+1),\qquad \eta=1/256.
\]
The candidate is supplied by the positive-domain recurrence in Section~1.
Induction checks its denotation; positivity verifies all denominators.
Its observed residual is zero, so choose $c_g=0$.

The spectral lower-bound hole can be discharged without approximate
numerical arithmetic. For all $u,v\in\R$,
\[
 \begin{pmatrix}u&v\end{pmatrix}
 (Y-I/16)\binom uv
 =\frac{15}{16}\left(u+\frac8{15}v\right)^2+\frac{v^2}{240}\geq0.
                                                               \tag{4}
\]
Thus $Y\succeq I/16$, so $\sigma_2(Y)\geq1/16$ and
$b_1(Y)=\sigma_2(Y)/2\geq1/32$. Choose $\beta=1/32$.
The remaining empirical hypothesis is exactly
\[
 h_{\mathrm{noise}}:\quad
 \norm{Y-\ev_c f_*}_F/2\leq1/256.                              \tag{5}
\]
Conditional on (5), the margin is positive and (3) gives
\[
 \forall A\in\mathrm{WFA}_{\leq1},\qquad
 L_{f_*}(A)-L_{f_*}(g)\geq
 \left(\frac7{256}\right)^2-\left(\frac1{256}\right)^2
 =\frac3{4096}.                                                \tag{6}
\]
The output of this elaboration is a proof of (5)$\Rightarrow$(6),
plus the candidate and its domain proof. For this export, (5) is kept
in the outer context or abstracted as an argument; it is not falsely
filled as a solved empirical hole. Without experimental or modeling
support for (5), it is not an unconditional empirical result.

This check also exposes what remains unresolved. The sequence equal to
$1,1/2,1/3$ at lengths $0,1,2$ and zero afterward has a three-state
nilpotent WFA realization and gives exactly the same $Y$.
Consequently the observed data do not canonically select the rational
recurrence over all finite WFAs. On the full infinite behavior, every
Hilbert Hankel block is positive definite because
\[
 v^TH_mv=\int_0^1\left(\sum_{i=0}^{m-1}v_it^i\right)^2dt>0
 \quad(v\ne0),
\]
so no finite WFA realizes $g$ exactly. This global proof depends on the
specified sequence, not on extrapolating the three observations.

\section{What canonicity and forgetting would mean here}

A canonicity field must name its admissible outputs and equivalence.
For example, minimal reachable and observable finite linear realizations
of a fixed full behavior are unique up to state-space isomorphism.
That theorem is available only when the full behavior has finite Hankel
rank. A canonical image realization of a repaired finite table is relative
to that table and its closure and consistency data. Neither statement
makes the proposed rational model a canonical scientific mechanism.
Likewise, uniqueness of an elaborated proof term, uniqueness of a model's
behavior, uniqueness of its state presentation, and minimum description
cost are different propositions and require different fibers of
$\Cat{Can}$.

The diagram's forgetful maps have identifiable losses.
\begin{enumerate}
\item Model to full behavior forgets implementation, hidden state,
      precision, and resource costs. Its kernel identifies observationally
      equivalent implementations.
\item Behavior to finite observations forgets all untested words. Its
      kernel is equality on the tested outputs, much coarser than global
      behavioral equality.
\item A finite observational image quotients prefix combinations by
      the kernel of their observation map. With fixed prefix generators,
      adding tests gives $K_{\mathrm{new}}\subseteq K_{\mathrm{old}}$
      and a canonical new-to-old quotient map. The identity on prefix
      generators descends in the old-to-new direction only if
      $K_{\mathrm{old}}\subseteq K_{\mathrm{new}}$, hence only if the
      kernels are equal in this setting. Arbitrary linear maps, such as
      the zero map, are not excluded; the obstruction is to this
      generator-compatible map.
\item Proof to proposition forgets the proof object. Further erasing
      its context would also lose assumptions and is not a sound
      promotion to an unconditional conclusion.
\item Language-tag erasure loses the class-relative resource measure.
      A coordinate count in one language cannot be compared as a universal
      complexity measure without a common cost model.
\end{enumerate}
The task objective decides which losses are acceptable. Prediction on a
fixed horizon can tolerate much more forgetting than mechanism discovery
under interventions. Such interventions are additional experimental
structure and are not introduced by a proof of word-behavior equality.

\section{Status and next obligations}

\paragraph{Proved mathematical fragments.}
The category-of-elements, commuting-square, and section constructions
give the displayed categories and their substitution/filling laws
conditional on the \emph{supplied} task presheaves, natural maps, and
contextual interface. The field schema is not a proof that this complete
encoding exists in a specific proof assistant. Exact semantic transports compose.
Approximate transports compose with explicit norm-dependent budgets.
The robust comparison theorem, rational recurrence, rational spectral
certificate, and finite-evidence countermodel have proofs above.

\paragraph{Kernel-checked fragment.}
The companion \texttt{ComparisonKernel.lean} checks the norm-gap inference,
the scalar squared-gap inference, additive approximation transport,
positivity and recurrence arithmetic for $1/(n+1)$, its reciprocal
coordinate equation, identity (4) and its quadratic lower bound, and the
rational value $3/4096$. These checks do not include the WFA/Hankel rank
factorization, low-rank approximation theorem, spectral interpretation
of (4), the categorical construction, or the empirical premise (5).
They are a small verified consumer of the typed premises, not a full
formalization of this note.

\paragraph{Specified but not implemented.}
The task signature as concrete substitution-stable presheaves, its natural
transformation packaging, the complete proof-producing elaborator,
extraction of its contextual syntax into this displayed diagram,
checker acceptance under substitution, tactic search and dependency tracking,
validated matrix certificate generation, and a sound statistical
noise-evidence producer remain to be built. None is implied by the
presence of a checked scalar inference rule.

\paragraph{A bounded useful next implementation.}
Implement the displayed task record and partial-solution telescope in
one chosen proof assistant; connect the existing matrix lemmas; accept
rational LDL certificates like (4); and make every exported comparison
retain its noise-event premise and exact language/norm/horizon.
A successful run should reject three specific category errors:
using a relaxed matrix as a realization, presenting a replacement as an
exact transport, and erasing an empirical assumption during export.
A canonicity claim should remain an unfilled goal unless its stated
uniqueness theorem is supplied.

This is one elaboration situation for this problem. Another scientific
setting may need intervention semantics, stochastic processes, or a
new comparison objective and a different diagram. The construction earns
its use here by making the actual proof dependencies and prohibited
compositions inspectable, rather than by asserting a universal form for
theory formation.

\begin{thebibliography}{9}
\bibitem{displayed}
B. Ahrens and P. L. Lumsdaine.
\emph{Displayed Categories}. Logical Methods in Computer Science
15(1:20), 2019.
\url{https://doi.org/10.23638/LMCS-15(1:20)2019}.
The source supplies the general displayed-category framework; the task
records and comparison instance here are a separate construction.
\bibitem{cwf}
P. Dybjer. \emph{Internal Type Theory}.
TYPES 1995, LNCS 1158, pp.~120--134, 1996.
\url{https://www.cse.chalmers.se/~peterd/papers/InternalTT.pdf}.
This supplies the category-with-families background used for contextual
syntax and substitution, not a pre-existing implementation of this task.
\end{thebibliography}
\end{document}
